We have introduced proportionality as a fair solution concept for centroid clustering. Although exact proportional solutions may not exist, we gave efficient algorithms for computing approximate proportional solutions, and considered constrained optimization and sampling for further applications. Finally, we studied proportionality on real data and observed a data dependent tradeoff between proportionality and the $k$-means objective. While this tradeoff is in some sense a negative result, it also demonstrates that proportionality as a fairness guarantee matters in the sense that it meaningfully constrains the space of solutions. 

We have shown that $\rho$-proportional solutions need not exist for $\rho < 2$, and always exist for $\rho \geq 1+\sqrt{2}$. Closing this approximability gap is one outstanding question. Another is whether there is a more efficient and easily interpretable algorithm for optimizing total cost subject to proportionality, as our approach in Section~\ref{sec:lp} requires solving a linear program on the entire data set. We would ideally like a more efficient and easily interpretable primal-dual or local search type algorithm. More generally, what other fair solution concepts for clustering should be considered alongside proportionality, and can we characterize their relative advantages and disadvantages? Finally, can the idea of proportionality as a group fairness concept be adapted for supervised learning tasks like classification and regression? 